\ifdefined\gkver\else\def\gkver{student}\fi
\documentclass[\gkver]{gaokaozhenti}
\begin{document}

\chapter{1999年全国}

\begin{enumerate}
\item
%% number: 1
%% paperName: 1999年全国高考第 1 题 4 分
%% typeId: 2
%% type: 多选题
%% chapter: 原子和原子核
%% point: 玻尔模型
%% method: 无
%% score: 4
%% degree: 750
%% duplicateId: 0
%% body:
下列说法正确的是\xzanswer{BC}
\fourchoices[answer=BC]
{当氢原子从 $n =$ 2 的状态跃迁到 $n =$ 6 的状态时，发射出光子}
{放射性元素的半衰期是指大量该元素的原子核中有半数发生衰变需要的时间}
{同一元素的两种同位素具有相同的质子数}
{中子与质子结合成氘核时吸收能量}
%% answer: BC
%% memo:
\memoanswer{
当氢原子核外电子从 $n=2$ 能级跃迁到 $n=6$ 能级时，是从低能级向高能级跃迁，需吸收光子，故 A 错误；B 选项叙述了半衰期的概念，故 B 正确；同位素是指具有相同质子数的一类原子，在元素周期表中占有相同位置，故 C 正确；中子与质子聚合是太阳中时刻进行的热核反应，放出大量热量，故 D 错误。
}

\item
%% number: 2
%% paperName: 1999年全国高考第 2 题 4 分
%% typeId: 1
%% type: 单选题
%% chapter: 电路
%% point: 闭合电路欧姆定律
%% method: 无
%% score: 4
%% degree: 850
%% duplicateId: 0
%% body:
一太阳能电池板，测得它的开路电压为 $800\UmV$，短路电流为 $40\UmA$。若将该电池板与一阻值为 $20\UO$ 的电阻器连成一闭合电路，则它的路端电压是\xzanswer{D}
\fourchoices[answer=D]
{$0.10\UV$}
{$0.20\UV$}
{$0.30\UV$}
{$0.40\UV$}
%% answer: D
%% memo:
\memoanswer{
由题意知 $E = 800\UmV$，短路电流为 $40\UmA$，由 $E = I_0 r$ 得 $r = \frac{E}{I_0} = \frac{800}{40} = 20\UO$，因外接电阻器的阻值与内阻相同，得路端电压为 $U = E \frac{R}{R + r} = 0.40\UV$，故 D 正确。
}

\item
%% number: 3
%% paperName: 1999年全国高考第 3 题 4 分
%% typeId: 2
%% type: 多选题
%% chapter: 原子和原子核
%% point: 天然放射性
%% method: 无
%% score: 4
%% degree: 850
%% duplicateId: 0
%% body:
下列说法正确的是\xzanswer{BC}
\fourchoices[answer=BC]
{\ce{^{226}_{88}Ra} 衰变为 \ce{^{222}_{86}Rn} 要经过 1 次 $\alpha$ 衰变和 1 次 $\beta$ 衰变}
{\ce{^{238}_{92}U} 衰变为 \ce{^{234}_{91}Pa} 要经过 1 次 $\alpha$ 衰变和 1 次 $\beta$ 衰变}
{\ce{^{232}_{90}Th} 衰变为 \ce{^{208}_{82}Pb} 要经过 6 次 $\alpha$ 衰变和 4 次 $\beta$ 衰变}
{\ce{^{238}_{92}U} 衰变为 \ce{^{222}_{86}Rn} 要经过 4 次 $\alpha$ 衰变和 4 次 $\beta$ 衰变}
%% answer: BC
%% memo:
\memoanswer{
核反应前后满足质量数和核电荷数守恒。设衰变前、后原子核质量数为 $M$、$M'$，核电荷数为 $Z$、$Z'$，$\alpha$ 衰变次数为 $m$，$\beta$ 衰变次数为 $n$，将四个选项衰变过程中质量数、核电荷数的变化分别代入 $M - M' = 4m$，$Z - Z' = 2m - n$ 中验证，故 BC 正确。
}

\item
%% number: 4
%% paperName: 1999年全国高考第 4 题 4 分
%% typeId: 2
%% type: 多选题
%% chapter: 气体性质
%% point: 理想气体状态方程
%% method: 无
%% score: 4
%% degree: 650
%% duplicateId: 0
%% body:
一定质量的理想气体处于平衡状态 Ⅰ。现设法使其温度降低而压强升高，达到平衡状态 Ⅱ，则\xzanswer{BC}
\fourchoices[answer=BC]
{状态 Ⅰ 时气体的密度比状态 Ⅱ 时的大}
{状态 Ⅰ 时分子的平均动能比状态 Ⅱ 时的大}
{状态 Ⅰ 时分子间的平均距离比状态 Ⅱ 时的大}
{状态 Ⅰ 时每个分子的动能都比状态 Ⅱ 时的分子平均动能大}
%% answer: BC
%% memo:
\memoanswer{
由一定质量理想气体的状态方程 $\frac{pV}{T}=C$ 知，$T$ 降低、$p$ 增大，$V$ 一定减小，气体密度 $\rho=\frac{m}{V}$，所以密度增大，故 A 错误 C 正确；因状态 Ⅰ 时温度高，所以分子的平均动能大，故 B 正确；分子的平均动能是对大量分子动能进行统计的结果，对单独一个分子的动能大小毫无意义，故 D 错误。
}

\item
%% number: 5
%% paperName: 1999年全国高考第 5 题 4 分
%% typeId: 1
%% type: 单选题
%% chapter: 光学
%% point: 光的折射
%% method: 无
%% score: 4
%% degree: 650
%% duplicateId: 0
%% body:
假设地球表面不存在大气层，那么人们观察到的日出时刻与实际存在大气层的情况相比\xzanswer{B}
\fourchoices[answer=B]
{将提前}
{将延后}
{在某些地区将提前，在另一些地区将延后}
{不变}
%% answer: B
%% memo:
\memoanswer{
如图所示，当光从真空射入空气时，发生向下的弯曲（如带箭头的实线），所以人们观察到的日出时刻比实际提前；假设地球表面不存在大气层，那么光线如图中虚线，则人们观察到的日出时刻与实际存在大气层的情况相比将延后，故 B 正确。

\onepicture[w=8cm]{figs/05_an.jpg}
}

\item
%% number: 6
%% paperName: 1999年全国高考第 6 题 4 分
%% typeId: 2
%% type: 多选题
%% chapter: 电磁感应
%% point: 右手定则
%% method: 无
%% score: 4
%% degree: 450
%% duplicateId: 0
%% body:
如图所示为地磁场磁感线的示意图。在北半球地磁场的竖直分量向下。飞机在我国上空匀速巡航，机翼保持水平，飞行高度不变。由于地磁场的作用，金属机翼上有电势差。设飞行员左方机翼末端处的电势为 $\varphi_{1}$，右方机翼末端处的电势为 $\varphi_{2}$，则\xzanswer{AC}
\onepicture[w=8cm]{figs/06.png}
\fourchoices[answer=AC]
{若飞机从西往东飞，$\varphi_{1}$ 比 $\varphi_{2}$ 高}
{若飞机从东往西飞，$\varphi_{2}$ 比 $\varphi_{1}$ 高}
{若飞机从南往北飞，$\varphi_{1}$ 比 $\varphi_{2}$ 高}
{若飞机从北往南飞，$\varphi_{2}$ 比 $\varphi_{1}$ 高}
%% answer: AC
%% memo:
\memoanswer{
如图所示，当飞机从西往东飞时，由右手定则可判断：感应电动势方向向北（如图所示），电势为左高右低，所以 $\varphi_{1}$ 比 $\varphi_{2}$ 高，当飞机从东往西时，$\varphi_{1}$ 仍比 $\varphi_{2}$ 高，故 A 正确 B 错误；当飞机从南往北飞时，由右手定则可判断：感应电动势方向向西（如图所示），电势为左高右低，所以 $\varphi_{1}$ 比 $\varphi_{2}$ 高，当飞机从北往南时，$\varphi_{1}$ 仍比 $\varphi_{2}$ 高，故 C 正确 D 错误。

\onepicture[w=11cm]{figs/06_an.jpg}
}

\item
%% number: 7
%% paperName: 1999年全国高考第 7 题 4 分
%% typeId: 1
%% type: 单选题
%% chapter: 交流电
%% point: 变压器
%% method: 无
%% score: 4
%% degree: 650
%% duplicateId: 0
%% body:
下面是 4 种亮度可调的台灯的电路示意图，它们所用的白炽灯泡相同，且都是“$220\UV$，$40\UW$”。当灯泡所消耗的功率都调至 $20\UW$ 时，哪种台灯消耗的功率最小？\xzanswer{C}
\fourchoices[answer=C, ispicture=true, h=3.6cm]
{figs/07a.jpg}
{figs/07b.jpg}
{figs/07c.jpg}
{figs/07d.jpg}
%% answer: C
%% memo:
\memoanswer{
台灯的消耗功率是指包含灯泡和其他辅助器材的总功率。C 中理想变压器功率损耗为零，电源输出的总功率（台灯消耗功率）只有灯泡的功率 $20\UW$，而其他选项中，不论滑动变阻器是分压接法（如 D）还是限流接法（如 A、B），滑动变阻器上总有功率损耗，台灯的消耗功率都大于 $20\UW$，故 C 正确。
}

\item
%% number: 8
%% paperName: 1999年全国高考第 8 题 4 分
%% typeId: 2
%% type: 多选题
%% chapter: 功和能
%% point: 动能定理
%% method: 无
%% score: 4
%% degree: 650
%% duplicateId: 0
%% body:
一物体静止在升降机的地板上，在升降机加速上升的过程中，地板对物体的支持力所做的功等于\xzanswer{CD}
\fourchoices[answer=CD]
{物体势能的增加量}
{物体动能的增加量}
{物体动能的增加量加上物体势能的增加量}
{物体动能的增加量加上克服重力所做的功}
%% answer: CD
%% memo:
\memoanswer{
在升降机加速上升的过程中，物体受重力 $mg$、地板施加的支持力 $N$ 的作用。重力对其做负功，支持力对其做正功。设升降机上升高度为 $h$，由动能定理得 $W_N - mgh = \Delta E_k$，$W_N = mgh + \Delta E_k$，其中 $mgh$ 为物体势能的增加量，也等于物体克服重力所做的功，$\Delta E_k$ 为物体动能的增加量，故 AB 错误 CD 正确。
}

\item
%% number: 9
%% paperName: 1999年全国高考第 9 题 4 分
%% typeId: 2
%% type: 多选题
%% chapter: 曲线运动
%% point: 竖直平面圆周运动
%% method: 无
%% score: 4
%% degree: 650
%% duplicateId: 0
%% body:
如图，细杆的一端与一小球相连，可绕过 O 点的水平轴自由转动。现给小球一初速度，使它做圆周运动，图中 a、b 分别表示小球轨道的最低点和最高点，则杆对球的作用力可能是\xzanswer{AB}
\onepicture[w=4.2cm]{figs/09.png}
\fourchoices[answer=AB]
{a 处为拉力，b 处为拉力}
{a 处为拉力，b 处为推力}
{a 处为推力，b 处为拉力}
{a 处为推力，b 处为推力}
%% answer: AB
%% memo:
\memoanswer{
小球在最低点（a 处）时所需向心力沿杆由 a 指向圆心 O，向心力是杆对小球的力与小球重力的合力，重力方向向下，故杆必定给球向上的力，所以 a 处一定是拉力；小球在最高点（b 处）时，若杆恰好对球没有作用力，即小球的重力恰好提供向心力，设此时小球速度为 $v_b$，则 $mg = m\frac{v_b^2}{R}$，当小球在最高点的速度 $v > v_b$，向心力大于重力，杆对小球有向下的拉力，以弥补重力提供向心力之不足；当小球在最高点的速度 $v < v_b$，向心力小于重力，杆对小球有向上的推力，以抵消重力提供向心力之多余，故 AB 正确。
}

\item
%% number: 10
%% paperName: 1999年全国高考第 10 题 4 分
%% typeId: 2
%% type: 多选题
%% chapter: 曲线运动
%% point: 人造地球卫星
%% method: 无
%% score: 4
%% degree: 450
%% duplicateId: 0
%% body:
地球同步卫星到地心的距离 $r$ 可由 $r^{3} = \frac{{{a^2}{b^2}c}}{{4{\pi ^2}}}$ 求出。已知式中 $a$ 的单位是 $\Um$，$b$ 的单位是 $\Us$，$c$ 的单位是 $\Umsq$，则\xzanswer{AD}
\fourchoices[answer=AD]
{$a$ 是地球半径，$b$ 是地球自转的周期，$c$ 是地球表面处的重力加速度}
{$a$ 是地球半径，$b$ 是同步卫星绕地心运动的周期，$c$ 是同步卫星的加速度}
{$a$ 是赤道周长，$b$ 是地球自转周期，$c$ 是同步卫星的加速度}
{$a$ 是地球半径，$b$ 是同步卫星绕地心运动的周期，$c$ 是地球表面处的重力加速度}
%% answer: AD
%% memo:
\memoanswer{
卫星与地球间的万有引力提供向心力，有 $G\frac{Mm}{r^{2}}=m\frac{4\pi^{2}}{T^{2}}r$，得 $r^{3}=\frac{GMT^{2}}{4\pi^{2}}$，又地球表面的重力加速度 $g=\frac{GM}{R^{2}}$，联立得 $r^{3}=\frac{R^{2}T^{2}g}{4\pi^{2}}$，式中 $R$ 是地球半径，$T$ 是地球自转周期，亦即同步卫星运行周期，$g$ 是地球表面处的重力加速度，对照式 $r^{3}=\frac{a^{2}b^{2}c}{4\pi^{2}}$，故 AD 正确。
}

\item
%% number: 11
%% paperName: 1999年全国高考第 11 题 4 分
%% typeId: 1
%% type: 单选题
%% chapter: 力物体的平衡
%% point: 物体系平衡
%% method: 整体与隔离
%% score: 4
%% degree: 450
%% duplicateId: 0
%% body:
如图所示，两木块的质量分别为 $m_{1}$ 和 $m_{2}$，两轻质弹簧的劲度系数分别为 $k_{1}$ 和 $k_{2}$，上面木块压在上面的弹簧上（但不拴接），整个系统处于平衡状态。现缓慢向上提上面的木块，直到它刚离开上面弹簧。在这过程中下面木块移动的距离为\xzanswer{C}
\onepicture[w=2.8cm]{figs/11.png}
\fourchoices[answer=C]
{$\frac{{{m_1}g}}{{{k_1}}}$}
{$\frac{{{m_2}g}}{{{k_1}}}$}
{$\frac{{{m_1}g}}{{{k_2}}}$}
{$\frac{{{m_2}g}}{{{k_2}}}$}
%% answer: C
%% memo:
\memoanswer{
整个系统处于平衡状态时，由胡克定律得，$k_2$ 弹簧的弹力为 $k_2 x_1 = (m_1 + m_2) g$；当 $m_1$ 被提离上面的弹簧时，$k_2$ 弹簧的弹力为 $k_2 x_2 = m_2 g$，故木块 $m_2$ 移动的距离为 $x_1 - x_2 = \frac{(m_1 + m_2) g}{k_2} - \frac{m_2 g}{k_2} = \frac{m_1 g}{k_2}$，故 C 正确。
}

\item
%% number: 12
%% paperName: 1999年全国高考第 12 题 4 分
%% typeId: 2
%% type: 多选题
%% chapter: 电磁感应
%% point: 电磁感应中图象
%% method: 无
%% score: 4
%% degree: 450
%% duplicateId: 0
%% body:
一匀强磁场，磁场方向垂直纸面，规定向里的方向为正。在磁场中有一细金属圆环，线圈平面位于纸面内，如图 \ref{1999全国12a} 所示。现令磁感强度 $B$ 随时间 $t$ 变化，先按图 \ref{1999全国12b} 中所示的 Oa 图线变化，后来又按图线 bc 和 cd 变化，令 $\varepsilon_{1}$、$\varepsilon_{2}$、$\varepsilon_{3}$ 分别表示这三段变化过程中感应电动势的大小，$I_{1}$、$I_{2}$、$I_{3}$ 分别表示对应的感应电流，则\xzanswer{BD}
\twopicture[labelA=1999全国12a, labelB=1999全国12b, wA=3.5cm, wB=7cm]
{figs/12a.png}
{figs/12b.png}
\fourchoices[answer=BD]
{$\varepsilon_{1} > \varepsilon_{2}$， $I_{1}$ 沿逆时针方向，$I_{2}$ 沿顺时针方向}
{$\varepsilon_{1} < \varepsilon_{2}$， $I_{1}$ 沿逆时针方向，$I_{2}$ 沿顺时针方向}
{$\varepsilon_{1} < \varepsilon_{2}$，$I_{2}$ 沿顺时针方向，$I_{3}$ 沿逆时针方向}
{$\varepsilon_{2} = \varepsilon_{3}$，$I_{2}$ 沿顺时针方向，$I_{3}$ 沿顺时针方向}
%% answer: BD
%% memo:
\memoanswer{
由图线 Oa 段 $B$ 大于零且随时间均匀增加判断该段时间内 $B$ 向里增强，且 $\frac{\Delta B}{\Delta t}=$ 恒量，由楞次定律判断感应电流 $I_1$ 方向为逆时针方向，由法拉第电磁感应定律得 $E_1 = \frac{\Delta \phi}{\Delta t} = \frac{\Delta B}{4} \cdot S$；对 bc 段，$B$ 大于零，即方向向里且逐渐减小，由楞次定律判断感应电流 $I_2$ 的方向为顺时针，由法拉第电磁感应定律得 $E_2 = \frac{\Delta \phi}{\Delta t} = \frac{\Delta B}{1} \cdot S$，式中 $\Delta B$ 与 $E_1$ 中 $\Delta B$ 相同，故 $E_2 = E_1$，即 $E_1 < E_2$；对 cd 段，$B$ 为负，方向向外，且逐渐增强，由楞次定律判断该段感应电流 $I_3$ 的方向为顺时针，由于 cd 段与 bc 段斜率相同，故 $\frac{\Delta B}{\Delta t}$ 相同，由法拉第电磁感应定律得 $E_3 = \frac{\Delta \phi}{\Delta t} = \frac{\Delta B}{1} \cdot S$，即 $E_3 = E_2$，综上 BD 正确。
}

\item
%% number: 13
%% paperName: 1999年全国高考第 13 题 5 分
%% typeId: 3
%% type: 填空题
%% chapter: 光学
%% point: 光的折射
%% method: 无
%% score: 5
%% degree: 850
%% duplicateId: 0
%% body:
一束单色光，在真空中波长为 $6.00\times10^{-7}\Um$，射入折射率为 1.50 的玻璃中。它在此玻璃中的波长是 \tkanswer[2cm]{$4.00\times10^{-7}\Um$}，频率是 \tkanswer[2.5cm]{$5.00\times10^{14}\UHz$}。（真空中光速是 $3.00\times10^{8}\Ums$）
%% answer: 4.00×10^{-7} m，5.00×10^{14} Hz
\jdanswer{
$4.00\times10^{-7}\Um$，$5.00\times10^{14}\UHz$
}
%% memo:
\memoanswer{
由 $n=\frac{c}{v}=\frac{\lambda_{0}}{\lambda}$，解得 $\lambda=\frac{\lambda_{0}}{n}=\frac{6.00\times10^{-7}}{1.50}\Um=4.00\times10^{-7}\Um$；由 $c=\lambda_{0}\nu$，解得 $\nu=\frac{c}{\lambda_{0}}=\frac{3.00\times10^{8}}{6.00\times10^{-7}}\UHz=5.00\times10^{14}\UHz$。
}

\item
%% number: 14
%% paperName: 1999年全国高考第 14 题 5 分
%% typeId: 3
%% type: 填空题
%% chapter: 直线运动
%% point: 竖直上下抛运动
%% method: 无
%% score: 5
%% degree: 650
%% duplicateId: 0
%% body:
一跳水运动员从离水面 $10\Um$ 高的平台上向上跃起，举双臂直体离开台面，此时其重心位于从手到脚全长的中点，跃起后重心升高 $0.45\Um$ 达到最高点。落水时身体竖直，手先入水。（在此过程中运动员水平方向的运动忽略不计。）从离开跳台到手触水面，他可用于完成空中动作的时间是 \tkanswer[1.5cm]{$1.7\Us$}。（计算时，可以把运动员看作全部质量集中在重心的一个质点。$g$ 取为 $10\Umsq$，结果保留二位数字。）
%% answer: 1.7 s（答 1.8 秒同样给分。）
\jdanswer{
$1.7\Us$（答 1.8 秒同样给分。）
}
%% memo:
\memoanswer{
上升阶段：运动员做竖直上抛运动，有 $t_1 = \sqrt{\frac{2h_1}{g}} = \sqrt{\frac{2 \times 0.45}{10}}\Us = 0.3\Us$，下落阶段：运动员做自由落体运动，有 $t_2 = \sqrt{\frac{2h_2}{g}} = \sqrt{\frac{2 \times 10}{10}}\Us = \sqrt{2}\Us = 1.4\Us$，所以 $t = t_1 + t_2 = 1.7\Us$。
}

\item
%% number: 15
%% paperName: 1999年全国高考第 15 题 5 分
%% typeId: 3
%% type: 填空题
%% chapter: 电场
%% point: 等势面
%% method: 无
%% score: 5
%% degree: 650
%% duplicateId: 0
%% body:
图中 A、B、C、D 是匀强电场中一正方形的四个顶点。已知 A、B、C 三点的电势分别为 $\varphi_{A} = 15\UV$，$\varphi_{B} = 3\UV$，$\varphi_{C} = -3\UV$。由此可得 D 点电势 $\varphi_{D} =$ \tkanswer[1.2cm]{$9\UV$}
\onepicture[w=4.5cm, align=right]{\includesvg{figs/15}}
%% answer: 9 V
\jdanswer{
$9\UV$
}
%% memo:
\memoanswer{
由匀强电场的特性知，等长的平行线段的两端电势差相同，由题意知 A、B 间电势差为 $12\UV$，B、C 间电势差为 $6\UV$，故 D、C 间电势差为 $12\UV$，A、D 间电势差 $6\UV$，故 D 点电势为 $9\UV$。
}

\item
%% number: 16
%% paperName: 1999年全国高考第 16 题 5 分
%% typeId: 3
%% type: 填空题
%% chapter: 机械波
%% point: 波的图象
%% method: 无
%% score: 5
%% degree: 450
%% duplicateId: 0
%% body:
图 \ref{1999全国16a} 中有一条均匀的绳，1、2、3、4、$\ldots$ 是绳上一系列等间隔的点。现有一列简谐横波沿此绳传播。某时刻，绳上 9、10、11、12 四点的位置和运动方向如图 \ref{1999全国16b} 所示（其他点的运动情况未画出），其中点 12 的位移为零，向上运动，点 9 的位移达到最大值。试在图 \ref{1999全国16c} 中画出再经过 $\frac{3}{4}$ 周期时点 3、4、5、6 的位置和速度方向，其他点不必画。（图 \ref{1999全国16c} 的横、纵坐标与图 \ref{1999全国16a}、\ref{1999全国16b} 完全相同）
\threepicture[labelA=1999全国16a, labelB=1999全国16b, labelC=1999全国16c, w=12cm, align=right]
{figs/16a.png}
{figs/16b.png}
{figs/16c.png}
%% answer: 如图
\jdanswer{
如图
}
%% memo:
\memoanswer{
题中给定 9，10，11，12 各质点某时刻的位置和运动方向，依次可以画出此时刻的波形图并知道这列波向右传播，进而明确 3，4，5，6 各质点此时的位置和运动方向。经过 $\frac{3}{4}T$ 后，这列波向右传播了 $\frac{3}{4}\lambda$，故将此时刻的波形图向右平移 $\frac{3}{4}\lambda$ 后得到了 $\frac{3}{4}T$ 后的波形图。由此波形图逐一弄清 3，4，5，6 各点的位置和速度方向，如图所示。

\onepicture[w=11cm]{figs/16_an.jpg}
}

\item
%% number: 17
%% paperName: 1999年全国高考第 17 题 4 分
%% typeId: 4
%% type: 实验题
%% chapter: 直线运动
%% point: 游标卡尺和螺旋测微仪
%% method: 无
%% score: 4
%% degree: 850
%% duplicateId: 0
%% body:
用游标为 50 分度的卡尺（测量值可准确到 $0.02\Umm$）测定某圆筒的内径时，卡尺上的示数如图。可读出圆筒的内径为 \tkanswer[1.5cm]{$52.12\Umm$}。
\onepicture[w=8cm, align=right]{figs/17.png}
%% answer: 52.12 mm
\jdanswer{
$52.12\Umm$
}
%% memo:
\memoanswer{
主尺读数为 $52\Umm$，游标每一分度是 $0.02\Umm$，其上第 6 条刻度线与主尺上毫米刻度线对齐，所以读数是 $52\Umm + 6 \times 0.02\Umm = 52.12\Umm$。
}

\item
%% number: 18
%% paperName: 1999年全国高考第 18 题 6 分
%% typeId: 4
%% type: 实验题
%% chapter: 光学
%% point: 透镜
%% method: 无
%% score: 6
%% degree: 650
%% duplicateId: 0
%% body:
某同学用小灯泡、凸透镜和光屏在水平光具座上做测量凸透镜焦距的实验。先使小灯泡、透镜和屏的中心等高，再调节三者间的距离，使屏上出现清晰的灯丝像。
\begin{enumerate}
	\item
	测量透镜到灯泡的距离 $l_{1}$ 及透镜到光屏的距离 $l_{2}$，则可得到透镜的焦距 $f =$ \tkanswer[2.5cm]{$\frac{l_{1}l_{2}}{l_{1}+l_{2}}$}（用 $l_{1}$，$l_{2}$ 表示）。
	\item
	保持透镜不动，将光屏向远离透镜的方向移动一个焦距的距离，再移动小灯泡，使屏上得到清晰的像。指出该像是放大还是缩小的，是正立还是倒立的，是实像还是虚像。答：\tkanswer[4cm]{放大倒立实像}。
\end{enumerate}
%% answer: （1）$f=\frac{l_1 l_2}{l_1+l_2}$（2）放大倒立实像。
\jdanswer{
\begin{enumerate}
	\item
	$f = \frac{l_1 l_2}{l_1 + l_2}$

	\item
	放大倒立实像。

\end{enumerate}
}
%% memo:
\memoanswer{
由透镜成像的公式 $\frac{1}{f}=\frac{1}{u}+\frac{1}{v}$ 得，解得 $f=\frac{l_{1}l_{2}}{l_{1}+l_{2}}$。

凸透镜成像规律：

\begin{itemize}
	\item
	物距小于焦距成正立放大虚像，应用是放大镜；
	\item
	物距大于一倍焦距小于二倍焦距成倒立放大实像，像距大于二倍焦距；
	\item
	物距等于二倍焦距成倒立等大实像，像距等于二倍焦距；
	\item
	物距大于二倍焦距成倒立缩小实像，像距大于一倍焦距小于二倍焦距。
\end{itemize}
}

\item
%% number: 19
%% paperName: 1999年全国高考第 19 题 7 分
%% typeId: 4
%% type: 实验题
%% chapter: 电路
%% point: 电路实验
%% method: 无
%% score: 7
%% degree: 450
%% duplicateId: 0
%% body:
图 \ref{1999全国19a} 为测量电阻的电路。$R_{x}$ 为待测电阻，$R$ 的阻值已知。$R'$ 为保护电阻，阻值未知。电源 $\varepsilon$ 的电动势未知。K$_{1}$、K$_{2}$ 均为单刀双掷开关。A 为电流表，其内阻不计。
\begin{enumerate}
	\item
	按图 \ref{1999全国19a} 所示的电路，在图 \ref{1999全国19b} 的实物图上连线。
	\item
	测量 $R_{x}$ 的步骤为：将 K$_{2}$ 向 d 闭合，K$_{1}$ 向 \tkanswer[0.8cm]{a} 闭合，记下电流表读数 $I_{1}$。将 K$_{2}$ 向 c 闭合，K$_{1}$ 向 \tkanswer[0.8cm]{b} 闭合，记电流表读数 $I_{2}$。计算 $R_{x}$ 的公式是 $R_{x} =$ \tkanswer[2.5cm]{$\frac{I_{2}R}{I_{1}}$}。
\end{enumerate}
\twopicture[labelA=1999全国19a, labelB=1999全国19b, wA=6cm, wB=7.5cm, align=right]
{figs/19a.png}
{figs/19b.png}
%% answer: （1）如图（2）a，b，$\frac{I_2 R}{I_1}$
\jdanswer{
\begin{enumerate}
	\item
	如图

	\item
	a，b，$\frac{I_2 R}{I_1}$

\end{enumerate}
}
%% memo:
\memoanswer{
将 $K_2$ 向 d 闭合，$K_1$ 向 a 闭合时，$R_x$ 与 $R$ 并联，$I_1$ 是流过 $R_x$ 的电流，当 $K_2$ 向 c 闭合，$K_1$ 向 b 闭合时，$R_x$ 与 $R$ 并联，$I_2$ 是流过 $R$ 的电流，$R'$ 是保护电阻，阻值应该很大，才能保证电路结构变化时起保护作用，也就是说 $R_x$ 和 $R$ 并联电路两端可以看做恒压，有 $I_1 R_x = I_2 R$，故 $R_x=\frac{I_2 R}{I_1}$。

\onepicture[w=7.5cm]{figs/19_an.jpg}
}

\item
%% number: 20
%% paperName: 1999年全国高考第 20 题 12 分
%% typeId: 6
%% type: 计算题
%% chapter: 运动和力
%% point: 动量守恒定律
%% method: 无
%% score: 12
%% degree: 650
%% duplicateId: 0
%% body:
试在下述简化情况下由牛顿定律导出动量守恒定律的表达式：系统是两个质点，相互作用力是恒力，不受其他力，沿直线运动。要求说明推导过程中每步的根据，以及式中各符号和最后结果中各项的意义。
%% answer: $m_1 v_1 + m_2 v_2 = m_1 v_1' + m_2 v_2'$
\jdanswer{
$m_{1}v_{1} + m_{2}v_{2} = m_{1}v_{1}' + m_{2}v_{2}'$
}
%% memo:
\memoanswer{
令 $m_1$ 和 $m_2$ 分别表示两质点的质量，$F_1$ 和 $F_2$ 分别表示它们所受的作用力，$a_1$ 和 $a_2$ 分别表示它们的加速度，$t_1$ 和 $t_2$ 分别表示 $F_1$ 和 $F_2$ 作用的时间。$v_1$ 和 $v_2$ 分别表示它们相互作用过程中的初速度，$v_1'$ 和 $v_2'$ 分别表示末速度，根据牛顿第二定律，有：

$F_1=m_1 a_1$，$F_2=m_2 a_2$ ①

由加速度的定义可知

$a_1=\frac{v_1' - v_1}{t_1}$，$a_2=\frac{v_2' - v_2}{t_2}$ ②

代入上式，可得

$F_1 t_1=m_1(v_1' - v_1)$，$F_2 t_2=m_2(v_2' - v_2)$ ③

根据牛顿第三定律，可知

$F_1=-F_2$；$t_1=t_2$ ④

由③，①可得

$m_1 v_1 + m_2 v_2 = m_1 v_1' + m_2 v_2'$ ⑤

其中 $m_1 v_1$ 和 $m_2 v_2$ 为两质点的初动量，$m_1 v_1'$ 和 $m_2 v_2'$ 为两质点的末动量，这就是动量守恒定律的表达式。

评分标准：本题 12 分。①、②、③各 1 分，④式 2 分，⑤式 3 分。正确、清楚说明每步的根据给 2 分，正确说出式中各符号和结果中各项意义的再给 2 分。
}

\item
%% number: 21
%% paperName: 1999年全国高考第 21 题 12 分
%% typeId: 6
%% type: 计算题
%% chapter: 直线运动
%% point: 多段匀变速直线运动
%% method: 无
%% score: 12
%% degree: 650
%% duplicateId: 0
%% body:
为了安全，在公路上行驶的汽车之间应保持必要的距离。已知某高速公路的最高限速 $v = 120\Ukmh$。假设前方车辆突然停止，后车司机从发现这一情况，经操纵刹车，到汽车开始减速所经历的时间（即反应时间）$t = 0.50\Us$。刹车时汽车受到的阻力的大小 $f$ 为汽车重力的 0.40 倍。该高速公路上汽车间的距离 $s$ 至少应为多少？$g$ 取 $10\Umsq$。
%% answer: 156 m
\jdanswer{
$156\Um$
}
%% memo:
\memoanswer{
在反应时间内，汽车作匀速运动，运动的距离

$s_1 = v t$ ①

设刹车时汽车的加速度的大小为 $a$，汽车的质量为 $m$，有

$f = m a$ ②

又 $f = 0.4 m g$

自刹车到停下，汽车运动的距离

$s_2 = \frac{v^2}{2a}$ ③

所求距离

$s = s_1 + s_2$ ④

由以上各式得

$s = 1.6\times10^{2}\Um$ ⑤

评分标准：本题 12 分。①、③、④式各 3 分，结果正确再给 3 分（结果为 $1.5\times10^{2}\Um$ 的，同样给分）
}

\item
%% number: 22
%% paperName: 1999年全国高考第 22 题 13 分
%% typeId: 6
%% type: 计算题
%% chapter: 电场
%% point: 带电粒子的偏转
%% method: 无
%% score: 13
%% degree: 450
%% duplicateId: 0
%% body:
在光滑水平面上有一质量 $m = 1.0\times10^{-3}\Ukg$、电量 $q = 1.0\times10^{-10}\UC$ 的带正电小球，静止在 $O$ 点。以 $O$ 点为原点，在该水平面内建立直角坐标系 $Oxy$。现突然加一沿 $x$ 轴正方向、场强大小 $E = 2.0\times10^{6}\UVm$ 的匀强电场，使小球开始运动。经过 $1.0\Us$，所加电场突然变为沿 $y$ 轴正方向，场强大小仍为 $E = 2.0\times10^{6}\UVm$ 的匀强电场。再经过 $1.0\Us$，所加电场又突然变为另一个匀强电场，使小球在此电场作用下经 $1.0\Us$ 速度变为零。求此电场的方向及速度变为零时小球的位置。
%% answer: 电场指向第三象限，与 $x$ 轴成 $225^\circ$ 角；$(0.40\Um, 0.20\Um)$
\jdanswer{
电场的方向指向第三象限，与 $x$ 轴成 $225\Udu$ 角，速度变为零时小球的位置为 $(0.40\Um, 0.20\Um)$。
}
%% memo:
\memoanswer{
由牛顿定律得知，在匀强电场中小球加速度的大小为

$qE = ma$，解得 $a = 0.20\Umsq$，①

当场强沿 $x$ 正方向时，经过 1 秒钟小球的速度大小为

$v_x = at = 0.20\times1.0 = 0.20\Ums$ ②

速度的方向沿 $x$ 轴正方向，小球沿 $x$ 轴方向移动的距离

$\Delta x_1 = \frac{1}{2}at^{2} = 0.10\Um$ ③，

在第 2 秒内，电场方向沿 $y$ 轴正方向，故小球在 $x$ 方向做速度为 $v_x$ 的匀速运动，在 $y$ 方向做初速为零的匀加速运动。沿 $x$ 方向移动的距离

$\Delta x_2 = v_x t = 0.20\Um$ ④

沿 $y$ 方向移动的距离

$\Delta y = \frac{1}{2}at^{2} = 0.10\Um$ ⑤，

故在第 2 秒末小球到达的位置坐标

$x_2 = \Delta x_1 + \Delta x_2 = 0.30\Um$ ⑥

$y_2 = \Delta y = 0.10\Um$ ⑦

在第 2 秒末小球在 $x$ 方向的分速度仍为 $v_x$，在 $y$ 方向的分速度

$v_y = at = 0.20\times1.0 = 0.20\Ums$ ⑧

由上可知，此时运动方向与 $x$ 轴成 $45\Udu$ 角。要使小球速度能变为零，则在第 3 秒内所加匀强电场的方向必须与此方向相反，即指向第三象限，与 $x$ 轴成 $225\Udu$ 角。

在第 3 秒内，设在电场作用下小球加速度的 $x$ 分量和 $y$ 分量分别为 $a_x$、$a_y$，则

$a_x = \frac{v_x}{t} = 0.20\Umsq$ ⑨，

$a_y = \frac{v_y}{t} = 0.20\Umsq$ ⑩，

在第 3 秒末小球到达的位置坐标为

$x_3 = x_2 + v_x t - \frac{1}{2}a_x t^{2} = 0.40\Um$ \textcircled{11}，

$y_3 = y_2 + v_y t - \frac{1}{2}a_y t^{2} = 0.20\Um$ \textcircled{12}。

小球的轨迹如图所示，小球的位置为 $(0.40\Um, 0.20\Um)$。

\onepicture[w=9cm]{figs/22_an.jpg}
}

\item
%% number: 23
%% paperName: 1999年全国高考第 23 题 13 分
%% typeId: 6
%% type: 计算题
%% chapter: 气体性质
%% point: 玻意耳定律
%% method: 无
%% score: 13
%% degree: 450
%% duplicateId: 0
%% body:
如图，气缸由两个横截面不同的圆筒连接而成。活塞 A、B 被轻质刚性细杆连接在一起，可无摩擦移动。A、B 的质量分别为 $m_{A} = 12\Ukg$，$m_{B} = 8.0\Ukg$，横截面积分别为 $S_{A} = 4.0\times10^{-2}\Umq$，$S_{B} = 2.0\times10^{-2}\Umq$。一定质量的理想气体被封闭在两活塞之间。活塞外侧大气压强 $p_{0} = 1.0\times10^{5}\UPa$。
\begin{enumerate}
	\item
	气缸水平放置达到如图 \ref{1999全国23a} 所示的平衡状态，求气体的压强；
	\item
	已知此时气体的体积 $V_{1} = 2.0\times10^{-2}\Umc$。现保持温度不变，将气缸竖直放置，达到平衡后如图 \ref{1999全国23b} 所示。与图 \ref{1999全国23a} 相比，活塞在气缸内移动的距离 $l$ 为多少？取重力加速度 $g = 10\Umsq$。
\end{enumerate}
\twopicture[labelA=1999全国23a, labelB=1999全国23b, hA=4cm, hB=4cm, align=right]
{figs/23a.png}
{figs/23b.png}
%% answer: （1）$p_1=1.0\times10^{5}$ Pa（2）$l=9.1\times10^{-2}$ m
\jdanswer{
\begin{enumerate}
	\item
	$p_1 = 1.0\times10^{5}\UPa$

	\item
	$l = 9.1\times10^{-2}\Um$

\end{enumerate}
}
%% memo:
\memoanswer{
（1）气缸处于图 1 位置时，设气缸内气体压强为 $p_1$，对于活塞和杆，力的平衡条件为

$p_0 S_A + p_1 S_B = p_1 S_A + p_0 S_B$ ①

解得 $p_1 = p_0 = 1.0\times10^{5}\UPa$ ②

（2）气缸处于图 2 位置时，设气缸内气体压强为 $p_2$，对于活塞和杆，力的平衡条件为

$p_0 S_A + p_2 S_B + (m_A + m_B)g = p_2 S_A + p_0 S_B$ ③

设 $V_2$ 为气缸处于图 2 位置时缸内气体的体积，由玻意耳定律可得

$p_1 V_1 = p_2 V_2$ ④

由几何关系可得

$V_1 - V_2 = l(S_A - S_B)$ ⑤

由以上各式解得

$l = 9.1\times10^{-2}\Um$ ⑥

评分标准：本题 13 分。第（1）问 4 分。①式 3 分，②式 1 分。第（2）问 9 分。③式 3 分，④、⑤、⑥式各 2 分。（最后结果在 $9.0\times10^{-2}\Um \sim 0.10\Um$ 之间的同样给分）
}

\item
%% number: 24
%% paperName: 1999年全国高考第 24 题 15 分
%% typeId: 6
%% type: 计算题
%% chapter: 磁场
%% point: 带电粒子在磁场中的运动
%% method: 无
%% score: 15
%% degree: 350
%% duplicateId: 0
%% body:
图中虚线 MN 是一垂直纸面的平面与纸面的交线，在平面右侧的半空间存在一磁感强度为 $B$ 的匀强磁场，方向垂直纸面向外。$O$ 是 MN 上的一点，从 $O$ 点可以向磁场区域发射电量为 $+q$、质量为 $m$、速率为 $v$ 的粒子，粒子射入磁场时的速度可在纸面内各个方向。已知先后射入的两个粒子恰好在磁场中给定的 $P$ 点相遇，$P$ 到 $O$ 的距离为 $L$。不计重力及粒子间的相互作用。
\begin{enumerate}
	\item
	求所考察的粒子在磁场中的轨道半径；
	\item
	求这两个粒子从 $O$ 点射入磁场的时间间隔。
\end{enumerate}
\onepicture[w=6.5cm, align=right]{figs/24.png}
%% answer: （1）$R=\frac{mv}{qB}$（2）$\Delta t=\frac{4m}{qB}\arccos\left(\frac{qLB}{2mv}\right)$
\jdanswer{
\begin{enumerate}
	\item
	$R = \frac{mv}{qB}$

	\item
	$\Delta t = \frac{4m}{qB}\arccos\left(\frac{qLB}{2mv}\right)$

\end{enumerate}
}
%% memo:
\memoanswer{
（1）设粒子在磁场中作圆周运动的轨道半径为 $R$，由牛顿第二定律，有

$qvB = m\frac{v^{2}}{R}$，

得 $R = \frac{mv}{qB}$ ①

（2）如图所示，以 OP 为弦可画两个半径相同的圆，分别表示在 P 点相遇的两个粒子的轨道。圆心和直径分别为 $O_1$、$O_2$ 和 $OO_1Q_1$、$OO_2Q_2$，在 O 处两个圆的切线分别表示两个粒子的射入方向，用 $\theta$ 表示它们之间的夹角。由几何关系可知

$\angle PO_1Q_1 = \angle PO_2Q_2 = \theta$ ②

从 O 点射入到相遇，粒子 1 的路程为半个圆周加弧长 $Q_1P$，由几何关系得

$P\hat{Q}_1 = R\theta$ ③，

粒子 2 的路程为半个圆周减弧长 $P\hat{Q}_2$，由几何关系得

$P\hat{Q}_2 = R\theta$ ④，

粒子 1 运动的时间

$t_1 = \frac{1}{2}T + \frac{R\theta}{v}$ ⑤，

其中 $T$ 为圆周运动的周期。粒子 2 运动的时间为

$t_2 = \frac{1}{2}T - \frac{R\theta}{v}$ ⑥，

两粒子射入的时间间隔

$\Delta t = t_1 - t_2 = 2\frac{R\theta}{v}$ ⑦，

因 $R\cos\frac{1}{2}\theta = \frac{1}{2}L$ 得 $\theta = 2\arccos\frac{L}{2R}$ ⑧，

由①、⑦、⑧三式得

$\Delta t = \frac{4m}{qB}\arccos\left(\frac{qLB}{2mv}\right)$ ⑨

\onepicture[w=6.5cm]{figs/24_an.jpg}

评分标准：本题 15 分。第（1）问 3 分。第（2）问 12 分。⑤、⑥式各 4 分，⑧、⑨式各 2 分。
}

\end{enumerate}

\end{document}
